% two_minds_1.tex --- first full draft
% Title alternatives (pick one):
%   Models of arithmetic of two minds about consistency
%   Of two minds: machine-cutoff provability and the intensionality of Con(PA)
%   Consistency verdicts are not model-invariant
% Authorship: drafted for Feldman--Hamkins; adjust as agreed.
\documentclass[11pt]{amsart}
\usepackage[margin=1.15in]{geometry}
\usepackage{amsmath,amssymb,amsthm}
\usepackage[colorlinks=true,linkcolor=blue!60!black,citecolor=blue!60!black,urlcolor=blue!60!black]{hyperref}
\usepackage{xcolor}

\newtheorem{theorem}{Theorem}[section]
\newtheorem{lemma}[theorem]{Lemma}
\newtheorem{corollary}[theorem]{Corollary}
\newtheorem{proposition}[theorem]{Proposition}
\theoremstyle{definition}
\newtheorem{definition}[theorem]{Definition}
\newtheorem{question}[theorem]{Question}
\theoremstyle{remark}
\newtheorem{remark}[theorem]{Remark}
\newtheorem{convention}[theorem]{Convention}

\newcommand{\PA}{\mathrm{PA}}
\newcommand{\Con}{\mathrm{Con}}
\newcommand{\Pf}{\mathrm{Proof}}
\newcommand{\Prov}{\mathrm{Pr}}
\newcommand{\HT}[1]{\mathrm{H}_{#1}}
\newcommand{\godel}[1]{\ulcorner #1 \urcorner}
\newcommand{\N}{\mathbb{N}}

% Draft-management macros
\newcommand{\verify}[1]{\textcolor{red!70!black}{\textbf{[verify:} #1\textbf{]}}}

\title[Of two minds about consistency]{Models of arithmetic of two minds about consistency}

\author{David Victor Feldman}
\address{Department of Mathematics and Statistics, University of New Hampshire}
\email{dvfinnh@gmail.com}

\date{August 2026}

\subjclass[2020]{Primary 03F40; Secondary 03F30, 03C62, 03H15, 68V20}

\keywords{Consistency statements, provability predicates, intensionality,
G\"odel's second incompleteness theorem, Rosser sentences, models of
arithmetic, formal verification}

\begin{document}

\begin{abstract}
For each $n$ we construct Turing machines $T_1,\dots,T_n$, none of which
halts, and associated provability predicates
$\Prov_{T_i}(x)$, asserting that $x$ has a proof appearing before $T_i$
halts. Each $\Prov_{T_i}$ is $\PA$-provably a subpredicate of the standard
provability predicate and defines exactly the standard provability relation
in the standard model; in particular $\PA$ proves
$\Con(\PA)\to\Con_{T_i}(\PA)$ for the associated consistency statements.
We show, assuming only the consistency of $\PA$, that these are the
\emph{only} implications: for every $S\subseteq\{1,\dots,n\}$ the theory
\[
\PA \;+\; \neg\Con(\PA) \;+\; \{\Con_{T_i}(\PA) : i\in S\}
        \;+\; \{\neg\Con_{T_i}(\PA) : i\notin S\}
\]
is consistent. In particular, a single model of $\PA$ can satisfy
$\Con_{T_1}(\PA)\wedge\neg\Con_{T_2}(\PA)$: whether a model of arithmetic
``thinks $\PA$ is consistent'' is not an invariant of the model, but
depends on the chosen arithmetization of provability. The proof combines
Rosser's witness-comparison method with the recursion theorem, relocating
the witness comparison from the syntax of a sentence into the dynamics of
a machine. We also prove a normal form theorem: for every machine $T$,
$\PA$ proves that $\Con_T(\PA)$ is equivalent to the disjunction of
$\Con(\PA)$ with the assertion that $T$ halts before any inconsistency
proof appears. Using it, we reduce the classification of the machine
consistency statements up to provable equivalence to two sharply posed
questions about $\Sigma_1$ sentences under witness comparison. The main
results have been formally verified in Lean~4.
\end{abstract}

\maketitle

\section{Introduction}

G\"odel's second incompleteness theorem asserts that a consistent,
sufficiently strong, effectively axiomatized theory such as $\PA$ neither
proves nor refutes its own consistency statement $\Con(\PA)$, and the
completeness theorem then yields models of $\PA$ satisfying $\Con(\PA)$
alongside models satisfying $\neg\Con(\PA)$. In a loose way of speaking,
one is tempted to say that some models of $\PA$ \emph{think} $\PA$ is
consistent while others think it is not, as though this were a property of
the model. But not so fast. The sentence $\Con(\PA)$ is not handed to us
by the theory; it is manufactured from choices---a G\"odel numbering, a
formalized proof predicate, an enumeration of the axioms---and the
independence phenomenon is notoriously sensitive to these choices. The
purpose of this note is to exhibit the sensitivity in a sharp form:
we produce a \emph{single} model of $\PA$ which is of two minds about the
consistency of $\PA$, satisfying one entirely natural-looking
arithmetization of ``$\,\PA$ is consistent'' while refuting another, where
both arithmetizations define exactly the true provability relation in the
standard model. More generally, we classify completely the possible
patterns of consistency verdicts for a natural family of such
arithmetizations.

That the consistency statement is intensional is of course a classical
insight. Rosser's theorem \cite{Rosser1936} already provides a one-line
witness: for the Rosser provability predicate
\[
\Prov^{R}(x) \;:=\; \exists y\,\bigl(\Pf(y,x)\wedge
   \forall z\le y\,\neg\Pf(z,\godel{0{=}1})\bigr),
\]
the associated consistency statement is an outright theorem of $\PA$, so
that any model of $\PA+\neg\Con(\PA)$ already satisfies
consistency-according-to-Rosser while refuting
consistency-according-to-Hilbert--Bernays---and the two predicates are
extensionally identical in the standard model (assuming, as we do
throughout, that $\PA$ is in fact consistent). Feferman's fundamental
study \cite{Feferman1960} showed that varying the \emph{enumeration of the
axioms} produces similar effects, and isolated conditions on an
arithmetization under which the second incompleteness theorem is
recovered; the derivability conditions of Hilbert--Bernays and L\"ob play
the same role on the side of the proof predicate. There is by now a rough
consensus among logicians that an acceptable formalization of provability
ought to be provably equivalent to the standard one, precisely because
predicates violating this permit the effects above; on this view the
deviant predicates do not \emph{genuinely express} provability, and the
second incompleteness theorem is not threatened. See Halbach and Visser
\cite{HalbachVisser2014a,HalbachVisser2014b} for a careful contemporary
analysis of the axes of intensionality, and \cite{Hamkins2022qeios} for
striking effects obtained purely on the axis of axiom enumeration.

We take no issue with that consensus; our point is that the consensus does
not dissolve the descriptive question but rather poses it. Grant that the
deviant predicates are deviant. They are nevertheless \emph{extensionally
correct}: in the real world---in the standard model---they define the very
same relation ``$y$ is a proof of $x$'' as the standard predicate, and the
ones constructed here are even $\PA$-provably \emph{sound} subpredicates
of the standard one, satisfying necessitation, and only ever
\emph{undercounting} proofs. How wildly can the consistency verdicts of
such predicates disagree, and how does the disagreement distribute across
the models of $\PA$? The answer, it turns out, is: as wildly as is not
outright provably impossible. We regard this as intensionality studied as
\emph{geography} rather than as pathology---a complete map of where the
verdicts can land---and the second incompleteness theorem itself supplies
the engine of the construction.

Our arithmetizations arise from a machine cutoff. For a Turing machine
$T$ which (in fact) never halts, let $\Prov_T(x)$ assert that $x$ has a
proof appearing strictly before $T$ halts, that is, a proof coded by some
$y$ such that $T$ has not halted within $y$ steps. Since $T$ never really
halts, $\Prov_T$ defines the standard provability relation in the standard
model; and provably in $\PA$, whatever $\Prov_T$ counts as a proof, so
does the standard predicate. In a model of $\PA+\neg\Con(\PA)$, however,
the verdict of the associated consistency statement $\Con_T(\PA)$ is
controlled by a race between two possibly nonstandard numbers: the least
proof of $0{=}1$, and the halting time of $T$, if any. The main theorem
states that this race can be rigged in every conceivable way,
simultaneously for any finite number of machines.

\begin{theorem}[Main Theorem; see Theorem~\ref{thm:main}]
\label{thm:main-intro}
Assume $\PA$ is consistent, and let $n\ge 1$. There exist Turing machines
$T_1,\dots,T_n$, none of which halts, such that:
\begin{enumerate}
\item each $\Prov_{T_i}$ defines the standard provability relation in the
standard model, and $\PA\vdash \Prov_{T_i}(x)\to\Prov(x)$, so in
particular $\PA\vdash\Con(\PA)\to\Con_{T_i}(\PA)$;
\item for every $S\subseteq\{1,\dots,n\}$, the theory
\[
\PA + \neg\Con(\PA) + \{\Con_{T_i}(\PA) : i\in S\}
    + \{\neg\Con_{T_i}(\PA) : i\notin S\}
\]
is consistent.
\end{enumerate}
Consequently the implications $\Con(\PA)\to\Con_{T_i}(\PA)$ generate
every provable implication among the sentences $\Con(\PA)$ and
$\Con_{T_1}(\PA),\dots,\Con_{T_n}(\PA)$. The verdict patterns
realizable in models of $\PA$ are then exactly these: all verdicts
true, or $\Con(\PA)$ false and the rest arbitrary.
\end{theorem}

Taking $n=2$ and $S=\{1\}$ yields the promised bifurcation.

\begin{corollary}\label{cor:two-minds-intro}
Assume $\PA$ is consistent. There are two arithmetizations of
provability, each defining the standard provability relation in the
standard model and each $\PA$-provably a subpredicate of the standard
predicate, together with a single model $M\models\PA$, such that $M$
satisfies the consistency of $\PA$ as formulated with the first and
refutes it as formulated with the second.
\end{corollary}

The proof is a Rosser argument, with the witness comparison relocated.
In Rosser's construction the witness comparison---``my proof comes before
any refutation''---is written into the \emph{syntax} of the sentence. Here
the comparison is relocated into the \emph{dynamics} of the machines: by
the recursion theorem, the machines $T_1,\dots,T_n$ jointly watch for a
$\PA$-proof that some verdict pattern is impossible, and upon finding one,
they conspire to \emph{instantiate} exactly that pattern, each machine
halting or entering an eternal loop according to its role in it. Any proof
constraining the pattern thereby defeats itself. No independence
hypotheses on the machines are needed; the construction discharges,
rather than assumes, any independence conditions (see Remark~\ref{rem:dynamics}). The technique
belongs to the witness-comparison tradition systematized by Guaspari and
Solovay \cite{GuaspariSolovay1979}, and the resulting consistency
statements $\Con_T(\PA)$ are $\Pi_1$ sentences strictly intermediate
between the provable and $\Con(\PA)$, cousins of the slow consistency
statement of Friedman, Rathjen and Weiermann \cite{FRW2013} and of
Pudl\'ak's partial consistency statements \cite{Pudlak1985}, with a
self-referentially timed rather than proof-theoretically graded cutoff.
We discuss these connections, and the precise fate of the derivability
conditions, in Section~\ref{sec:remarks}. In Section~\ref{sec:which} we
turn from constructing particular machine consistency statements to
classifying all of them: a normal form theorem
(Proposition~\ref{prop:normalform}) identifies the sentences $\Con_T(\PA)$,
up to $\PA$-provable equivalence, with the disjunctions of $\Con(\PA)$
with sentences asserting a Rosser-style victory of a machine over
inconsistency, and the classification problem then resolves into two
questions about $\Sigma_1$ sentences under witness comparison, which we
pose but do not settle.

Everything below assumes only the consistency of $\PA$; no soundness or
$\omega$-consistency assumptions are used,\footnote{With one scoping
caveat, which the formalization mentioned below makes precise: clauses
(1)--(3) of Theorem~\ref{thm:main} use only the consistency of $\PA$
(this is machine-verified). The final classification clause additionally
uses the consistency of $\PA+\Con(\PA)$ --- equivalently
$\PA\nvdash\neg\Con(\PA)$ --- which follows from the arithmetical
soundness of $\PA$ but not from its consistency alone.} and nothing is
special about $\PA$ beyond what is needed for the usual coding
(Remark~\ref{rem:base-theory}).

The main results of this paper --- Theorem~\ref{thm:main} in full
(clauses (1)--(3), including the necessitation and
extensional-correctness parts of clause (1)),
Corollary~\ref{cor:two-minds}, Proposition~\ref{prop:normalform}, and
Proposition~\ref{prop:realization} in the stage-formula form of
Remark~\ref{rem:pacing} --- have been formally verified in Lean~4 over
the Foundation library, with no unproved obligations and the standard
axiom footprint. The verified development takes the purely syntactic
form described in Remark~\ref{rem:syntactic}, and is carried out over an
arbitrary $\Delta_1$-axiomatized recursively enumerable extension of
$\mathsf{I}\Sigma_1$; the specialization of the hypotheses to $\PA$
itself additionally invokes the library's stated-but-unproved instance
of the standard fact that $\PA$ is $\Delta_1$-axiomatized. The
verification artifact is available at
\texttt{https://github.com/DavidVFeldman/two-minds-consistency}
and archived at \texttt{https://doi.org/10.5281/zenodo.21881585}.

\section{Machine-cutoff provability predicates}\label{sec:defs}

\begin{convention}\label{conv:coding}
Fix a standard G\"odel numbering and a standard proof predicate
$\Pf(y,x)$ for $\PA$, with $\Prov(x):=\exists y\,\Pf(y,x)$ and
$\Con(\PA):=\neg\Prov(\godel{0{=}1})$, satisfying the usual properties:
$\Pf$ is $\Delta_1$ in $\PA$, each proof codes a sequence of formulas
with a uniquely determined conclusion (a convenience only; see the
tie-break in the proof of Theorem~\ref{thm:main}), and the
Hilbert--Bernays--L\"ob derivability conditions hold for $\Prov$. For a
Turing machine $T$ (with a fixed standard coding of machines and
computations), let $\HT{T}(y)$ denote the formula ``$T$, run on empty
input, halts within $y$ steps.'' We arrange the coding so that, provably
in $\PA$: $\HT{T}(y)$ is $\Delta_1$; $\HT{T}$ is monotone, i.e.\
$\HT{T}(y)\wedge y\le z\to\HT{T}(z)$; and computations are provably
deterministic. For standard numerals these are the usual facts about
formalized finite simulations. In fact all that is ever used of
$\HT{T}$ below is that it is a provably monotone $\Delta_1$ predicate of
$y$; the machine, the determinism, and the halting-state bookkeeping are
presentational, a point the formalization exploits
(Remark~\ref{rem:syntactic}). % [bookkeeping: all standard; in a weak-base-theory version one would take H_T as Delta_0(exp) and track induction, cf. Remark on base theory]
\end{convention}

\begin{definition}\label{def:PrT}
For a Turing machine $T$, define
\[
\Prov_T(x) \;:=\; \exists y\,\bigl(\Pf(y,x)\wedge\neg\HT{T}(y)\bigr),
\qquad
\Con_T(\PA) \;:=\; \neg\Prov_T(\godel{0{=}1}).
\]
Thus $\Prov_T$ counts exactly those proofs which appear before $T$ halts,
and all proofs if $T$ never halts. Note that $\Con_T(\PA)$ is
$\PA$-provably equivalent to the $\Pi_1$ sentence
$\forall y\,(\Pf(y,\godel{0{=}1})\to\HT{T}(y))$.
\end{definition}

\begin{lemma}\label{lem:sub}
For every machine $T$:
\begin{enumerate}
\item $\PA\vdash\forall x\,(\Prov_T(x)\to\Prov(x))$; hence
      $\PA\vdash\Con(\PA)\to\Con_T(\PA)$.
\item If $T$ never halts, then $\Prov_T$ and $\Prov$ define the same
      relation in the standard model $\N$, namely true provability in
      $\PA$; moreover, for every standard sentence $\varphi$, if
      $\PA\vdash\varphi$ then $\PA\vdash\Prov_T(\godel{\varphi})$
      (necessitation).
\end{enumerate}
\end{lemma}

\begin{proof}
(1) is immediate from the definition. For (2): if $T$ never halts, then
for each standard $y$ the sentence $\neg\HT{T}(\bar y)$ is a true
$\Delta_1$ statement about a finite simulation, hence provable in $\PA$;
so in $\N$ the clause $\neg\HT{T}(y)$ is vacuously satisfied and
$\Prov_T$ coincides with $\Prov$, while if $\PA\vdash\varphi$ with proof
coded by the standard number $y$, then
$\PA\vdash\Pf(\bar y,\godel{\varphi})\wedge\neg\HT{T}(\bar y)$, giving
$\PA\vdash\Prov_T(\godel{\varphi})$.
\end{proof}

Lemma~\ref{lem:sub}(1) already settles half of the classification: in any
model of $\PA+\Con(\PA)$, every verdict $\Con_T(\PA)$ comes out true.
All the freedom, if there is any, lives inside the models of
$\PA+\neg\Con(\PA)$, where it is governed by a race between the least
proof of $0{=}1$ and the halting of $T$. The next section makes the race
explicit.

\section{Pattern sentences}\label{sec:patterns}

Fix machines $T_1,\dots,T_n$. For a set $S\subseteq\{1,\dots,n\}$, define
the \emph{pattern sentence}
\[
\sigma_S \;:=\; \exists p\,\Bigl(\Pf(p,\godel{0{=}1})
   \;\wedge\; \forall q<p\,\neg\Pf(q,\godel{0{=}1})
   \;\wedge\; \bigwedge_{i\in S}\HT{T_i}(p)
   \;\wedge\; \bigwedge_{i\notin S}\neg\HT{T_i}(p)\Bigr),
\]
asserting that there is a least proof $p$ of $0{=}1$, and that exactly the
machines $T_i$ with $i\in S$ have halted by stage $p$.

\begin{lemma}[Reformulation Lemma]\label{lem:pattern}
For each $S\subseteq\{1,\dots,n\}$,
\[
\PA\;\vdash\;\sigma_S\;\longleftrightarrow\;
\Bigl(\neg\Con(\PA)\wedge\bigwedge_{i\in S}\Con_{T_i}(\PA)
      \wedge\bigwedge_{i\notin S}\neg\Con_{T_i}(\PA)\Bigr).
\]
Consequently, the theory
$\PA+\neg\Con(\PA)+\{\Con_{T_i}(\PA):i\in S\}
+\{\neg\Con_{T_i}(\PA):i\notin S\}$
is consistent if and only if $\PA\nvdash\neg\sigma_S$.
\end{lemma}

\begin{proof}
Work in $\PA$. Suppose $\sigma_S$, witnessed by the least inconsistency
proof $p$. Certainly $\neg\Con(\PA)$. For $i\in S$: given any $y$ with
$\Pf(y,\godel{0{=}1})$, minimality gives $p\le y$, and monotonicity of
$\HT{T_i}$ gives $\HT{T_i}(y)$; as $y$ was arbitrary, $\Con_{T_i}(\PA)$
holds in the $\Pi_1$ form of Definition~\ref{def:PrT}. For $i\notin S$:
the witness $p$ itself satisfies
$\Pf(p,\godel{0{=}1})\wedge\neg\HT{T_i}(p)$, so
$\Prov_{T_i}(\godel{0{=}1})$, i.e.\ $\neg\Con_{T_i}(\PA)$.

Conversely, suppose $\neg\Con(\PA)$ holds together with the displayed
pattern of verdicts. By the least number principle there is a
least inconsistency proof $p$. For $i\in S$, applying $\Con_{T_i}(\PA)$
to $p$ yields $\HT{T_i}(p)$. For $i\notin S$, $\neg\Con_{T_i}(\PA)$ gives
some $y$ such that $\Pf(y,\godel{0{=}1})$ and $\neg\HT{T_i}(y)$;
minimality gives $p\le y$, and monotonicity (contrapositively) gives
$\neg\HT{T_i}(p)$. So $p$ witnesses $\sigma_S$.

The final claim is the completeness theorem plus the observation that the
displayed theory is $\PA$-provably equivalent to $\PA+\sigma_S$.
\end{proof}

The pattern sentences are, we would suggest, the real objects of study
here: the classification problem for verdict patterns is exactly the
question of which $\neg\sigma_S$ are provable, and the machines will now
be built so that none of them is.

\section{The construction and the main theorem}\label{sec:main}

\begin{theorem}\label{thm:main}
Assume $\PA$ is consistent, and let $n\ge1$. There exist Turing machines
$T_1,\dots,T_n$, none of which halts, such that
$\PA\nvdash\neg\sigma_S$ for every $S\subseteq\{1,\dots,n\}$.
Consequently, by Lemmas~\ref{lem:sub} and~\ref{lem:pattern}:
\begin{enumerate}
\item each $\Prov_{T_i}$ defines true $\PA$-provability in $\N$, is
$\PA$-provably a subpredicate of $\Prov$, and satisfies necessitation;
\item $\PA\vdash\Con(\PA)\to\Con_{T_i}(\PA)$ for each $i$; and
\item for every $S\subseteq\{1,\dots,n\}$ the theory
\[
\PA+\neg\Con(\PA)+\{\Con_{T_i}(\PA):i\in S\}
   +\{\neg\Con_{T_i}(\PA):i\notin S\}
\]
is consistent.
\end{enumerate}
In particular, the only provable implications among
$\Con(\PA),\Con_{T_1}(\PA),\dots,\Con_{T_n}(\PA)$ are those generated by
the implications $\Con(\PA)\to\Con_{T_i}(\PA)$.
\end{theorem}

\begin{proof}
\emph{The construction.} By the recursion theorem (in its standard
$n$-fold form, or by applying the ordinary recursion theorem to a single
controller index from which the $n$ machines are uniformly computed), we
may construct machines $T_1,\dots,T_n$, each with access to the full
tuple of indices, behaving as follows. Each $T_i$ performs the same
synchronized search, scanning $y=0,1,2,\dots$ in order and checking
whether $y$ codes a $\PA$-proof whose conclusion is one of the $2^n$
\emph{watched sentences} $\neg\sigma_S$, $S\subseteq\{1,\dots,n\}$
(these sentences are computed from the indices, and are pairwise
distinct; if a single stage codes proofs of several watched sentences,
the trigger adopts the least pattern among them, in some fixed ordering
of the patterns --- this tie-break makes the construction independent of
any uniqueness convention for conclusions). If the search finds a least
trigger $y_0$, with adopted pattern $S^*$, then:
\begin{itemize}
\item if $i\in S^*$, machine $T_i$ halts immediately;
\item if $i\notin S^*$, machine $T_i$ enters a designated absorbing loop
state, in which it remains forever.
\end{itemize}
If the search never triggers, the machines search forever. Thus the
family reads the first metamathematical verdict about itself and
conspires to instantiate exactly the pattern just ``refuted.'' We arrange
the loop state explicitly so that $\PA$ proves, by induction, that a
machine in the loop state at some stage is in it at every later stage
and never halts.

\emph{No watched sentence is provable.} Suppose toward a contradiction
that $\PA\vdash\neg\sigma_{S}$ for some $S$; then the search really
triggers, say least at $y_0$ (a standard number), with adopted pattern
$S^*$ (which need not be the $S$ we started from --- the contradiction
below runs through $S^*$). Note that $\neg\sigma_{S^*}$ is really
provable, $y_0$ coding a proof of it. The subsequent
behavior of each machine is a concrete finite computational fact: there
is a standard number $t$, computable from $y_0$, such that each $T_i$
with $i\in S^*$ halts within $t$ steps, and each $T_i$ with $i\notin S^*$
enters the loop state within $t$ steps. Being true $\Delta_1$ statements
about finite simulations, these facts are provable:
\[
\PA\vdash\HT{T_i}(\bar t)\ \ (i\in S^*),
\qquad
\PA\vdash\forall y\,\neg\HT{T_i}(y)\ \ (i\notin S^*),
\]
the latter by combining the provable finite fact that $T_i$ is in the
loop state at stage $\bar t$ with the provable absorption of the loop
state.

Now reason inside $\PA+\neg\Con(\PA)$. By the least number principle
there is a least inconsistency proof $p$. Suppose first $p>\bar t$. Then
for $i\in S^*$, monotonicity and $\HT{T_i}(\bar t)$ give $\HT{T_i}(p)$;
and for $i\notin S^*$ we have $\neg\HT{T_i}(p)$ outright. So $p$
witnesses $\sigma_{S^*}$---contradicting the $\PA$-provable
$\neg\sigma_{S^*}$. Hence $\PA+\neg\Con(\PA)\vdash p\le\bar t$, and so
\[
\PA+\neg\Con(\PA)\;\vdash\;\exists q\le\bar t\,\Pf(q,\godel{0{=}1}).
\]
But $t$ is standard and $\PA$ is (really) consistent, so for each
standard $q\le t$ the sentence $\neg\Pf(\bar q,\godel{0{=}1})$ is a true
$\Delta_1$ sentence, hence $\PA$-provable; conjoining these finitely many
facts, $\PA\vdash\forall q\le\bar t\,\neg\Pf(q,\godel{0{=}1})$.
Therefore $\PA+\neg\Con(\PA)$ is inconsistent, so
$\PA\vdash\Con(\PA)$---contradicting the second incompleteness theorem,
which applies since $\PA$ is consistent.

\emph{Conclusion.} No watched sentence is provable; hence in the real
world the search never triggers, no $T_i$ ever halts, and
Lemma~\ref{lem:sub} yields (1) and (2). For (3), since
$\PA\nvdash\neg\sigma_S$, the theory $\PA+\sigma_S$ is consistent, and
Lemma~\ref{lem:pattern} identifies it with the displayed theory. For the
final claim: any purported provable implication not generated by
$\Con(\PA)\to\Con_{T_i}(\PA)$ would exclude one of the verdict patterns
shown consistent in (2) and (3), or a pattern with all verdicts true,
which holds in any model of $\PA+\Con(\PA)$.
\end{proof}

\begin{corollary}[A model of two minds]\label{cor:two-minds}
Assume $\PA$ is consistent. With $n=2$ and $S=\{1\}$ in
Theorem~\ref{thm:main}, there is a model $M\models\PA$ with
\[
M\models \Con_{T_1}(\PA)\wedge\neg\Con_{T_2}(\PA),
\]
although $\Prov_{T_1}$ and $\Prov_{T_2}$ both define true
$\PA$-provability in the standard model, both are $\PA$-provably sound
subpredicates of the standard provability predicate, and both satisfy
necessitation. Whether a model of $\PA$ ``thinks $\PA$ is consistent'' is
thus not a property of the model alone.
\end{corollary}

\begin{remark}[The syntactic reading]\label{rem:syntactic}
The construction admits an equivalent, purely syntactic description,
which is the form taken by the machine-checked proof. By the
simultaneous diagonal lemma there are sentences $\delta_S$
($S\subseteq\{1,\dots,n\}$), pairwise distinct, with
$\PA\vdash\delta_S\leftrightarrow\neg\sigma_S$, where the stage
predicates standing in for the $\HT{T_i}$ are explicit, non-self-referential
$\Delta_1$ formulas reading the codes $\godel{\delta_{S'}}$ as fixed
numerals: ``some stage $y_0\le y$ is least among stages coding a proof
of some $\delta_{S'}$, and the least pattern proved at $y_0$ contains
$i$.'' The witness comparison, which the machine formulation relocates
from syntax into dynamics, is thereby relocated back: the family
$(\delta_S)$ is a simultaneous Rosser-style fixed point, each member
asserting that its own pattern is not the one instantiated at the least
trigger, and no recursion theorem beyond the diagonal lemma is needed.
The two descriptions are intertranslatable; the syntactic one is the
shorter to verify formally.
\end{remark}

\begin{remark}[Nonstandard dynamics; independence discharged]
\label{rem:dynamics}
Boolean independence of the halting statements of the machines, and their
joint independence from $\Con(\PA)$, would not suffice for
Theorem~\ref{thm:main}: those hypotheses do not control the \emph{relative
position} of nonstandard witnesses, and what the pattern sentences $\sigma_S$
assert is precisely such a position, namely that $T_i$ halts below the least
inconsistency proof. The construction supplies the positional control by
self-reference rather than by independence: any proof constraining the verdict
pattern defeats itself.

Inside a model
$M\models\PA+\neg\Con(\PA)$. Since every sentence is $M$-provable there,
each watched sentence has a (nonstandard) proof in $M$, so from $M$'s
internal point of view the search \emph{does} trigger, at some
nonstandard stage, and the machines then halt or loop according to the
$M$-least trigger's pattern $S^*_M$. The theorem says that the model may
be chosen so that this internal drama plays out in any prescribed
position relative to the least inconsistency proof. In the real world,
meanwhile, the search runs forever and nothing ever happens: the machines
are, one might say, waiting for a contradiction that never comes.
\end{remark}

\section{Remarks, connections, and questions}\label{sec:remarks}

\begin{remark}[The derivability conditions, and how the escape works]
\label{rem:HBL}
Each $\Prov_{T_i}$ satisfies the first derivability condition
(necessitation, Lemma~\ref{lem:sub}), but must fail the second: a
$\PA$-verified closure of $\Prov_{T_i}$ under modus ponens would, with
necessitation and provable $\Sigma_1$-completeness for $\Prov_{T_i}$,
put $\Prov_{T_i}$ back under the hypotheses of the second incompleteness
theorem in L\"ob's form, and the analysis below shows $\Con_{T_i}(\PA)$
behaves incompatibly with that. Concretely the failure is visible:
modus ponens applied to proofs below the cutoff produces a longer proof,
which the cutoff may decapitate. The third condition (provable
$\Sigma_1$-completeness \emph{for $\Prov_{T_i}$ itself}) fails for the
same reason: certifying a $\Prov_{T_i}$-fact requires a proof appearing
in time. It is worth contrasting the shape of the escape with Rosser's:
Rosser's consistency statement becomes an outright theorem of $\PA$,
whereas our $\Con_{T_i}(\PA)$ remains unprovable---indeed
$\PA+\neg\Con_{T_i}(\PA)$ is consistent by Theorem~\ref{thm:main}---and
unrefutable, even over $\PA+\neg\Con(\PA)$. The deviation from the
standard predicate manifests not as a collapse of the consistency
statement into provability but as its detachment from $\neg\Con(\PA)$:
the standard verdict no longer decides the deviant one.
\end{remark}

\begin{remark}[Intermediate consistency statements]\label{rem:intermediate}
Each $\Con_{T_i}(\PA)$ is a $\Pi_1$ sentence with
$\PA\vdash\Con(\PA)\to\Con_{T_i}(\PA)$,
$\PA\nvdash\Con_{T_i}(\PA)$, and
$\PA+\Con_{T_i}(\PA)\nvdash\Con(\PA)$: a nontrivial consistency-like
statement strictly intermediate between the provable and full
$\Con(\PA)$. In spirit, $\Con_T(\PA)$ is a partial or cut-off consistency
statement---``there is no inconsistency proof of size below the halting
time of $T$''---placing it in the company of Pudl\'ak's finite and
restricted consistency statements \cite{Pudlak1985} and of the slow
consistency of Friedman, Rathjen, and Weiermann \cite{FRW2013} (see also
Henk and Pakhomov \cite{HenkPakhomov}), but with a cutoff timed
self-referentially by the recursion theorem rather than graded by a
proof-theoretic hierarchy. It remains to locate the
$\Con_{T}(\PA)$ of Theorem~\ref{thm:main} in the interpretability degrees
alongside those statements. % [tone down / expand after we check FRW and Pudlak carefully]
\end{remark}

\begin{remark}[Other axes of intensionality]\label{rem:axes}
The construction here varies the proof predicate while holding the
G\"odel numbering and axiom enumeration fixed. Feferman
\cite{Feferman1960} showed how much can be achieved by varying the
enumeration of axioms alone; a modern instance of that trick, in which a
$\PA$-enumeration is arranged to acquire the consistency strength of
large cardinal hypotheses, appears in \cite{Hamkins2022qeios}. Halbach
and Visser \cite{HalbachVisser2014a,HalbachVisser2014b} taxonomize the
axes systematically. It seems fair to describe the present contribution,
against that background, as: the machine-cutoff formulation; the
complete classification of verdict patterns for finitely many
extensionally correct, provably sound subpredicates; and the resulting
single-model bifurcation phenomenon of
Corollary~\ref{cor:two-minds}.
\end{remark}

\begin{remark}[Base theory]\label{rem:base-theory}
We have stated everything for $\PA$ for concreteness, but the argument
uses only: $\Delta_1$ representation of the relevant predicates with
provable monotonicity, the least number principle for $\Delta_1$
formulas, provable $\Sigma_1$-completeness for the standard predicate at
standard instances, the diagonal lemma, and the second incompleteness
theorem. Accordingly, the machine-checked development mentioned in the
introduction is carried out over an arbitrary $\Delta_1$-axiomatized
recursively enumerable theory extending $\mathsf{I}\Sigma_1$, and
Theorem~\ref{thm:main}, Corollary~\ref{cor:two-minds}, and
Proposition~\ref{prop:normalform} hold at that generality; the
statements for $\PA$ are instances.
\end{remark}

Further questions appear at the end of the next section.

\section{Which sentences are machine consistency statements?}
\label{sec:which}

Question: which sentences $\theta$ arise, up to $\PA$-provable
equivalence, as $\Con_T(\PA)$ for a machine $T$ that never halts? By
Lemma~\ref{lem:sub}, $\PA\vdash\Con(\PA)\to\theta$ is necessary, and one
might naively conjecture it sufficient, which would make the machine
consistency statements a normal form for all $\Pi_1$ consequences of
consistency. The situation is otherwise. The key is to make explicit the Rosser-style race that has
been implicit throughout.

\begin{definition}\label{def:HaltR}
For a machine $T$, let $\mathrm{HaltAt}_T(s)$ be the $\Delta_1$ formula
asserting that $T$ halts exactly at step $s$ (provably, such $s$ is
unique if it exists), and define
\[
\mathrm{Halt}^R(T) \;:=\; \exists s\,\bigl(\mathrm{HaltAt}_T(s)
   \wedge \forall q<s\,\neg\Pf(q,\godel{0{=}1})\bigr),
\]
a $\Sigma_1$ sentence asserting that \emph{$T$ halts before any
inconsistency proof appears}---that $T$ wins the Rosser race against
inconsistency. For a never-halting $T$, $\mathrm{Halt}^R(T)$ is false
in $\N$.
\end{definition}

\begin{proposition}[Normal form]\label{prop:normalform}
For \emph{every} machine $T$,
\[
\PA\;\vdash\;\Con_T(\PA)\;\longleftrightarrow\;
\Con(\PA)\vee\mathrm{Halt}^R(T).
\]
Equivalently, with $\sigma_{\{1\}}$ the pattern sentence of
Section~\ref{sec:patterns} for the single machine $T$:
$\PA\vdash\sigma_{\{1\}}\leftrightarrow
\neg\Con(\PA)\wedge\mathrm{Halt}^R(T)$.
\end{proposition}

\begin{proof}
Work in $\PA$. ($\leftarrow$) $\Con(\PA)\to\Con_T(\PA)$ is
Lemma~\ref{lem:sub}. Suppose $\mathrm{Halt}^R(T)$, witnessed by $s$, and
let $y$ be any inconsistency proof; then $y\ge s$ by the witness clause,
so $\HT{T}(y)$ by monotonicity, giving $\Con_T(\PA)$ in its $\Pi_1$ form.
($\rightarrow$) Suppose $\Con_T(\PA)$; if $\Con(\PA)$ we are done, so
suppose also $\neg\Con(\PA)$, with least inconsistency proof $p$.
Applying $\Con_T(\PA)$ to $p$ gives $\HT{T}(p)$, so $T$ halts at some
$s\le p$; and for $q<s\le p$, minimality of $p$ gives
$\neg\Pf(q,\godel{0{=}1})$. Thus $\mathrm{Halt}^R(T)$.
\end{proof}

Thus every machine consistency statement asserts exactly:
\emph{either $\PA$ is consistent, or the machine beat the inconsistency
to the tape.} The classification problem therefore splits along the two
disjuncts. Writing $[\,\cdot\,]$ for $\PA$-provable equivalence classes,
Proposition~\ref{prop:normalform} gives
\[
\{[\Con_T(\PA)] : T\text{ never halts}\}
\;=\;
\{[\Con(\PA)\vee\rho] : \rho\in\mathcal{H}\},
\]
where
\[
\mathcal{H}:=\{[\mathrm{Halt}^R(T)] : T\text{ never halts}\},
\]
and $\mathcal{H}$ consists of ($\equiv$-classes of) $\Sigma_1$ sentences
false in $\N$. The next proposition shows $\mathcal{H}$ is plentiful:
every $\Sigma_1$ sentence false in $\N$ contributes its
\emph{Rosser truncation} to $\mathcal{H}$, at a pacing determined by a
machine implementation.

\begin{proposition}[Realization]\label{prop:realization}
Let $\rho=\exists x\,\delta(x)$ be a $\Sigma_1$ sentence
($\delta\in\Delta_0$, without loss of generality) false in $\N$. Let
$T_\rho$ be the machine that checks $\delta(0),\delta(1),\dots$ in order
and halts upon first finding a witness, and let $c(x)$ denote the
$\PA$-provably total, provably strictly monotone function giving the step
at which $T_\rho$ completes the check of $x$. Then $T_\rho$ never halts,
and
\[
\PA\;\vdash\;\mathrm{Halt}^R(T_\rho)\;\longleftrightarrow\;\rho^{R,c},
\qquad\text{where}\quad
\rho^{R,c}:=\exists x\,\bigl(\delta(x)\wedge
   \forall q<c(x)\,\neg\Pf(q,\godel{0{=}1})\bigr).
\]
Consequently $\Con(\PA)\vee\rho^{R,c}$ is, up to provable equivalence, a
machine consistency statement, namely $\Con_{T_\rho}(\PA)$.
\end{proposition}

\begin{proof}
$T_\rho$ never halts since $\rho$ is false in $\N$. Work in $\PA$. If
$\mathrm{Halt}^R(T_\rho)$ holds with witness $s$, then provably $s=c(x_0)$
for the least witness $x_0$ of $\delta$, and the clause
$\forall q<s\,\neg\Pf(q,\godel{0{=}1})$ is exactly the truncation clause
at $x_0$; so $\rho^{R,c}$ holds. Conversely, given $x$ with $\delta(x)$
and no inconsistency proof below $c(x)$, the least witness $x_0\le x$
satisfies $c(x_0)\le c(x)$ by monotonicity, so $T_\rho$ halts at
$c(x_0)$ with no inconsistency proof below it, giving
$\mathrm{Halt}^R(T_\rho)$.
\end{proof}

\begin{remark}[The pacing is an artifact]\label{rem:pacing}
The pacing function $c$ belongs to the machine presentation, not to the
mathematics. Working directly with monotone $\Delta_1$ stage formulas
$\eta(y)$ in place of machines --- as the formalization does
(Remark~\ref{rem:syntactic}) --- one may take
$\eta_\rho(y):=\exists x\le y\,\delta(x)$, and the first stage at which
$\eta_\rho$ fires \emph{is} the least witness of $\rho$: the truncation
of Proposition~\ref{prop:realization} lands at $c=\mathrm{id}$, i.e.,
at
\[
\rho^{R}\;:=\;\exists x\,\bigl(\delta(x)\wedge\forall x'<x\,\neg\delta(x')
\wedge\forall q<x\,\neg\Pf(q,\godel{0{=}1})\bigr).
\]
Question~\ref{q:obstructionB} is accordingly better phrased
presentation-theoretically: is every $\Sigma_1$ sentence false in $\N$
provably equivalent to $\pi^{R}$ for \emph{some} $\Sigma_1$ presentation
$\pi$ of it? The sensitivity of $\pi^R$ to the choice of presentation is
then precisely the Guaspari--Solovay phenomenon. The stage-formula form
of Proposition~\ref{prop:realization} just described --- with
$\eta_\rho$ in place of $T_\rho$ and truncation at the least witness ---
is among the machine-verified results.
\end{remark}

\begin{corollary}[Reduction of the classification]\label{cor:reduction}
A sentence $\theta$ is $\PA$-provably equivalent to $\Con_T(\PA)$ for
some never-halting $T$ if and only if
$\PA\vdash\theta\leftrightarrow\Con(\PA)\vee\rho$ for some
$\rho\in\mathcal{H}$. In particular, the naive conjecture above holds if
and only if both of the following do:
\begin{enumerate}
\item[\textup{(A)}] every $\Pi_1$ sentence $\theta$ with
$\PA\vdash\Con(\PA)\to\theta$ satisfies: $\theta\wedge\neg\Con(\PA)$ is
$\PA$-provably equivalent to a $\Sigma_1$ sentence; and
\item[\textup{(B)}] every $\Sigma_1$ sentence false in $\N$ is
$\PA$-provably equivalent to $\mathrm{Halt}^R(T)$ for some never-halting
$T$---equivalently, is provably equivalent to a Rosser truncation
$\rho^{R,c}$ of some $\Sigma_1$ sentence at some machine pacing $c$.
\end{enumerate}
\end{corollary}

\begin{proof}
The displayed equality of classes gives the first claim. For the second:
if $\PA\vdash\Con(\PA)\to\theta$, then provably
$\theta\leftrightarrow(\theta\wedge\Con(\PA))\vee(\theta\wedge\neg\Con(\PA))
\leftrightarrow\Con(\PA)\vee(\theta\wedge\neg\Con(\PA))$, so $\theta$ is
of the required form if and only if $\theta\wedge\neg\Con(\PA)$ is
equivalent to a member of $\mathcal{H}$; condition (A) supplies
$\Sigma_1$-ness and condition (B) upgrades $\Sigma_1$-ness (falsity in
$\N$ being automatic, since $\theta\wedge\neg\Con(\PA)$ is false in $\N$
by the actual consistency of $\PA$) to membership in $\mathcal{H}$.
\end{proof}

We suspect both (A) and (B) fail, so that the machine consistency
statements form a proper---and, by
Proposition~\ref{prop:realization}, still remarkably rich---subclass of
the $\Pi_1$ consequences of consistency; but we have not proved either
failure, and we record the two halves as questions.

\begin{question}[Obstruction A]\label{q:obstructionA}
Is there a $\Pi_1$ sentence $\theta$ with
$\PA\vdash\Con(\PA)\to\theta$ such that $\theta\wedge\neg\Con(\PA)$ is
not $\PA$-provably equivalent to any $\Sigma_1$ sentence? The slow
consistency statement of \cite{FRW2013} seems a natural candidate: over
$\PA+\neg\Con(\PA)$ with least inconsistency proof $p$, slow consistency
asserts a $\Pi_1$-flavored domination condition tying the convergence of
$F_{\varepsilon_0}$ to the position of $p$, with no evident $\Sigma_1$
reformulation. Preservation of $\Sigma_1$ sentences under end-extensions
suggests a strategy: exhibit a model of
$\PA+\neg\Con(\PA)+\mathrm{Con}^{\mathrm{slow}}$ with an end-extension
to a model of $\PA+\neg\Con(\PA)+\neg\mathrm{Con}^{\mathrm{slow}}$.
% [The end-extension must preserve p as least inconsistency proof ---
%  automatic, since initial segments and Sigma_1 facts are preserved and
%  no smaller proofs appear. What must be arranged is new convergence of
%  F_{eps_0} at a point with insufficient room below p. Check against
%  FRW's model constructions; this may be close to their interpolation
%  arguments.]
\end{question}

\begin{question}[Obstruction B]\label{q:obstructionB}
Is every $\Sigma_1$ sentence false in $\N$ provably equivalent to
$\mathrm{Halt}^R(T)$ for some never-halting $T$ --- equivalently
(Remark~\ref{rem:pacing}), to a Rosser truncation $\pi^R$ of some
$\Sigma_1$ presentation $\pi$ of it? The truncation changes the
sentence, in general:
$\pi^{R}$ provably implies $\pi$, but the converse asks every model's
witness to appear before its inconsistency proofs, which is a
genuine positional demand. The case
$\PA\vdash\rho\to\neg\Con(\PA)$ is the hardest: there,
$\mathrm{Halt}^R(T)\equiv\rho$ would require the machine's Rosser
victory---which the machine achieves precisely by seeing \emph{no}
inconsistency proof before halting---to provably certify that
inconsistency proofs exist beyond its halting time. We expect the
witness-comparison methods of Guaspari and Solovay
\cite{GuaspariSolovay1979} (see also Lindstr\"om
\cite{Lindstrom2003}) to be the right tools here, in either direction.
% [Even the Boolean closure of the class H is unclear, for the same
%  timing reasons: running two machines in parallel halves the clock,
%  which shifts the truncation point. Worth a lemma either way.]
\end{question}

\begin{question}\label{q:infinite}
Theorem~\ref{thm:main} handles each finite $n$, with machines depending
on $n$. Is there a single uniformly computable infinite family
$(T_i)_{i\in\omega}$, none halting, realizing every ``locally
permitted'' verdict pattern---for instance, such that for every set
$A\subseteq\omega$ that is $\Sigma_1$-definable in a given model of
$\PA+\neg\Con(\PA)$\dots\ ---in an appropriate formulation? Already the
right statement of the infinitary classification is not evident, since
$\sigma_S$ becomes infinitary.
\end{question}

\begin{question}\label{q:iterate}
The map $T\mapsto\Con_T(\PA)$ invites iteration: one may form the
sentence $\Con_T(\PA+\Con_{T'}(\PA))$, and so on. What hierarchy
results, under provability or under interpretability, and how does it
interact with the classical progressions of iterated consistency?
\end{question}

\subsection*{Acknowledgments}
The author is grateful to Joel David Hamkins for generous advice and
correspondence on these questions, in particular concerning the
intensionality of provability predicates and the effects obtainable by
varying the enumeration of axioms. The formal verification reported in
the introduction was carried out with Aristotle (Harmonic), over the
Foundation library of the Formalized Formal Logic project.
% [add grant info if any]

\begin{thebibliography}{99}

\bibitem{Feferman1960}
S.~Feferman,
\emph{Arithmetization of metamathematics in a general setting},
Fund. Math. \textbf{49} (1960), 35--92.

\bibitem{FRW2013}
S.-D.~Friedman, M.~Rathjen, and A.~Weiermann,
\emph{Slow consistency},
Ann. Pure Appl. Logic \textbf{164} (2013), no.~3, 382--393.

\bibitem{GuaspariSolovay1979}
D.~Guaspari and R.~M.~Solovay,
\emph{Rosser sentences},
Ann. Math. Logic \textbf{16} (1979), no.~1, 81--99.

\bibitem{HalbachVisser2014a}
V.~Halbach and A.~Visser,
\emph{Self-reference in arithmetic I},
Rev. Symb. Log. \textbf{7} (2014), no.~4, 671--691.

\bibitem{HalbachVisser2014b}
V.~Halbach and A.~Visser,
\emph{Self-reference in arithmetic II},
Rev. Symb. Log. \textbf{7} (2014), no.~4, 692--712.

\bibitem{Hamkins2022qeios}
J.~D.~Hamkins,
\emph{Nonlinearity in the hierarchy of large cardinal consistency
strength}, Qeios (2022), \url{https://www.qeios.com/read/T63YDP}.
% [verify exact title/year against the posted version]

\bibitem{HenkPakhomov}
P.~Henk and F.~Pakhomov,
\emph{Slow and ordinary provability for Peano arithmetic},
preprint, arXiv:1602.01822.
% [verify status --- may have appeared; update]

\bibitem{Lindstrom2003}
P.~Lindstr\"om,
\emph{Aspects of Incompleteness},
Lecture Notes in Logic \textbf{10}, 2nd ed., A K Peters, 2003.

\bibitem{Pudlak1985}
P.~Pudl\'ak,
\emph{Cuts, consistency statements and interpretations},
J. Symbolic Logic \textbf{50} (1985), no.~2, 423--441.

\bibitem{Rosser1936}
J.~B.~Rosser,
\emph{Extensions of some theorems of G\"odel and Church},
J. Symbolic Logic \textbf{1} (1936), no.~3, 87--91.

\end{thebibliography}

\end{document}
